%% -*- coding: utf-8 -*-
%% Elsevier-compatible LaTeX article
%% Ceará Digital Twin for Wildfire Monitoring

\documentclass[preprint,12pt,times,numbers]{elsarticle}

\makeatletter
\gdef\ttdefault{cmtt}
\makeatother

\usepackage[english]{babel}
\usepackage[utf8]{inputenc}
\usepackage[T1]{fontenc}
\usepackage{amsmath,amssymb}
\usepackage{graphicx}
\usepackage{hyperref}
\newcommand{\orcidlink}[1]{\href{https://orcid.org/#1}{\textcolor{green!60!black}{\scriptsize ORCID}}}
\usepackage{booktabs}
\usepackage{xcolor}
\usepackage{xspace}
%\usepackage{microtype}
\usepackage{microtype}
\usepackage{listings}
\usepackage{caption}
\usepackage{subcaption}
\usepackage{float}

% Float placement — discourage float-only pages, keep text flow
\setcounter{topnumber}{2}
\setcounter{bottomnumber}{2}
\setcounter{totalnumber}{4}
\setcounter{dbltopnumber}{2}
\renewcommand{\topfraction}{0.85}
\renewcommand{\bottomfraction}{0.70}
\renewcommand{\textfraction}{0.20}
\renewcommand{\floatpagefraction}{0.70}
\setkeys{Gin}{width=\linewidth,keepaspectratio}

%% Unified table typography — use resulttable for all manuscript tables.
\makeatletter
\newenvironment{resulttable}[2]{% #1=caption, #2=label
  \begin{table}[!htbp]
  \centering
  \caption{#1}
  \label{#2}
  \small
  \setlength{\tabcolsep}{5pt}
  \renewcommand{\arraystretch}{1.08}
}{%
  \end{table}
}
\newcommand{\paneltitle}[1]{\textbf{#1}\par\vspace{0.35em}}
\makeatother


% Inline figures — always at source position; use width <= \linewidth
\newcommand{\figurainline}[4][\linewidth]{%
  \par\medskip\noindent
  \begin{minipage}{\linewidth}
    \centering
    \includegraphics[width=#1]{#2}%
    \captionof{figure}{#3}%
    \label{#4}%
  \end{minipage}
  \par\medskip
}

\definecolor{codegreen}{rgb}{0,0.6,0}
\definecolor{codegray}{rgb}{0.5,0.5,0.5}
\definecolor{codepurple}{rgb}{0.58,0,0.82}
\definecolor{backcolour}{rgb}{0.95,0.95,0.92}

\lstdefinestyle{mystyle}{
    backgroundcolor=\color{backcolour},
    commentstyle=\color{codegreen},
    keywordstyle=\color{magenta},
    numberstyle=\tiny\color{codegray},
    stringstyle=\color{codepurple},
    basicstyle=\ttfamily\footnotesize,
    breakatwhitespace=false,
    breaklines=true,
    captionpos=b,
    keepspaces=true,
    numbers=left,
    numbersep=5pt,
    showspaces=false,
    showstringspaces=false,
    showtabs=false,
    tabsize=2
}
\lstset{style=mystyle}

\journal{Environmental Modelling \& Software}
\title{Ceará Digital Twin for Wildfires: An Open-Source Agentic AI Platform Integrating Near Real-Time Satellite Data and LangGraph Orchestration}

%% Authors
\author[unifor,cor1]{Nauber Gois\corref{cor1}}
\ead{naubergois@gmail.com}
\author[usj]{Alexandre Lobo}
\ead{alexandrelobo@unifor.br}
\author[unifor]{Vladia Celia Monteiro Pinheiro}
\ead{vladiacelia@unifor.br}
\author[saxion]{Daniel Valente de Macedo}
\ead{d.valentedemacedo@saxion.nl}
\author[unifor2]{Adriano Bessa Albuquerque}
\ead{adrianoba@unifor.br}

\address[unifor]{Universidade de Fortaleza (UNIFOR), Programa de P\'{o}s-Gradua\c{c}\~{a}o em Inform\'{a}tica Aplicada, Fortaleza, CE, Brazil}
\address[usj]{Universidade de S\~{a}o Jos\'{e} (USJ), Macau, China}
\cortext[cor1]{Corresponding author at Universidade de Fortaleza (UNIFOR). E-mail: \texttt{naubergois@gmail.com}, \texttt{francisco.gois@cge.ce.gov.br}. Author ORCID identifiers are registered in the Editorial Manager submission form.}
\address[saxion]{Saxion University of Applied Sciences, Deventer, Netherlands}
\address[unifor2]{Universidade de Fortaleza (UNIFOR), Programa de P\'{o}s-Gradua\c{c}\~{a}o em Inform\'{a}tica Aplicada, Fortaleza, CE, Brazil}

\begin{document}

\begin{frontmatter}

\begin{abstract}
Wildfire monitoring in Brazil depends on centralized systems with substantial detection-to-availability latency. We present the \emph{Ceará Wildfire Monitoring Platform}, an open-source agentic system for environmental decision support. The platform integrates NASA FIRMS and Open-Meteo into a LangGraph pipeline (10~nodes) for ingestion, validation, parallel analysis, evidence fusion, risk classification, ReAct diagnosis, and three-class alerting. On INPE validation during the 2025 dry season (184~days, 24{,}573 hotspots), alerts achieved 84.2\% precision and 91.8\% recall, with bootstrap 95\% confidence intervals for YES-class precision. GOES-16 supports research benchmarks only (grid-scale F1$\leq$0.05 vs.\ INPE). NeKo-PIGNN attains $R^2=0.972$ on synthetic dynamics; on short real-data windows, simpler baselines yield lower RMSE. The system deploys on a single EC2 instance (approximately USD~20/month) under the MIT license.
\end{abstract}

\begin{keyword}
Digital Twin \sep Wildfires \sep Physics-Informed GNN \sep Koopman Operator \sep Three-Class Detection \sep LangGraph \sep NASA FIRMS
\end{keyword}

\end{frontmatter}

%% ===================================================================
\section{Introduction}
\label{sec:intro}
%% ===================================================================

Wildfires represent one of the greatest environmental threats in Brazil, with direct impacts on public health, biodiversity, and the economy \cite{article:alencar2020,master:silva2021}. The State of Ceará, located in the northeastern semi-arid region, exhibits strong drought seasonality and sensitive biomes such as the Caatinga, making continuous monitoring an operational necessity \cite{techreport:funceme2024}.

Traditional monitoring systems such as INPE BDQueimadas and NASA FIRMS provide essential data but exhibit three operational gaps: 3--6~hour latency between satellite overpass and data availability \cite{manual:firms2024}; no automatic fusion with local climate and municipal context; and interfaces aimed at remote sensing specialists rather than civil defense managers.

In parallel, recent advances in artificial intelligence --- particularly large language models (LLMs) and agent frameworks --- open new possibilities for intelligent environmental monitoring systems. The paradigm of agents with ReAct reasoning \cite{article:react2023} and state graph orchestration (LangGraph \cite{manual:langgraph2024}) enables building auditable \emph{pipelines} that combine heterogeneous sources and produce natural language explanations.

We adopt \emph{Digital Twin} terminology pragmatically: the platform continuously mirrors satellite and weather observations and hosts predictive modules, but it does not yet provide full bidirectional control or physics-coupled simulation at operational scale~\cite{article:bauer2023natrev}. Throughout the paper we refer to the deployable artifact as the \emph{monitoring platform}.

This paper contributes:

\begin{enumerate}
    \item \textbf{An open-source platform} that integrates NASA FIRMS (VIIRS/MODIS), Open-Meteo, and Nominatim into a single pipeline, with optional GOES-16 for research benchmarks, without relying on an institutional database;
    \item \textbf{A LangGraph-orchestrated agent pipeline} that performs collection $\rightarrow$ Pydantic validation $\rightarrow$ parallel analysis (geospatial, GOES-16, climate) $\rightarrow$ evidence fusion $\rightarrow$ risk classification $\rightarrow$ ReAct diagnosis $\rightarrow$ alerts;
    \item \textbf{An RAG module} with FAISS index and local embeddings that enables natural language queries about the system's methodology, architecture, and usage;
    \item \textbf{Operational results} with real data collected between July 2024 and June 2026, demonstrating the feasibility of monitoring based exclusively on open sources.
\end{enumerate}

%% ===================================================================
\section{Related Work}
\label{sec:related}
%% ===================================================================

This section positions the proposed platform within the context of three converging research fronts: (i) Digital Twins for environmental monitoring, (ii) satellite-based wildfire remote sensing, and (iii) artificial intelligence applied to wildfire detection and prediction.

\subsection{Environmental Digital Twins}

The Digital Twin paradigm, originally conceived for manufacturing where a virtual model mirrors and controls a physical system in real time~\cite{article:grieves2014,article:tao2019}, has been systematically surveyed across definitions, architectures, and challenges~\cite{article:barricelli2019,article:rasheed2020,article:fuller2020,article:jones2020}. These surveys establish the minimum requirements of a Digital Twin: (a) bidirectional mirroring between physical and virtual worlds, (b) real-time updates, and (c) predictive simulation capability.

The expansion of the paradigm to environmental and Earth systems gained traction with the European Destination Earth~\cite{article:voosen2020} and proposals for Earth System Digital Twins~\cite{article:bauer2021nature,article:bauer2023natrev,article:nativi2021}. Bauer et al.~\cite{article:bauer2023natrev} argue that environmental Digital Twins differ fundamentally from industrial ones by operating at much larger spatiotemporal scales, with inherent uncertainties in natural processes and multi-source integration requirements.

\subsection{Digital Twins for Wildfires}

Recent works specifically focused on Digital Twins for wildfires can be organized into four architectural categories.

\paragraph{Agent-oriented approaches.} IVSR (Intelligent Virtual Situation Room)~\cite{article:morsali2026} is the most conceptually aligned proposal with this work. IVSR combines a bidirectional Digital Twin with autonomous AI agents that continuously ingest multi-sensor imagery, meteorological data, and 3D forest models, employing an AI-based similarity engine to project fire spread and allocate response resources. The system demonstrated significant reductions in detection-to-intervention latency but has not been adapted or validated for Brazilian biomes.

\paragraph{Vision-Language and RL approaches.} FIRE-VLM~\cite{article:webb2026} introduces an innovative framework combining a CLIP-style Vision-Language Model with reinforcement learning (PPO) to guide UAVs in tracking wildfires within a physics-grounded Digital Twin. Results show up to 6$\times$ reduction in detection time across five evaluation tasks. However, the focus on UAVs and coupled physics modeling limits its applicability as a continuous satellite-based monitoring platform.

\paragraph{Multi-modal sensing approaches.} FIRETWIN~\cite{article:raha2025} combines Unreal Engine with the CAWFE fire model for immersive reconstruction of historical wildfires (King Fire), integrating multi-sensor data and interactive analysis. AIMNET~\cite{article:aimnet2025} proposes a three-layer IoT-empowered Digital Twin (physical world $\rightarrow$ bidirectional links $\rightarrow$ digital world) for continuous gas emission monitoring with distributed sensor networks for early carbon plume detection. H-RDT~\cite{article:hrdt2025} adopts PINNs with multimodal fusion to model wildfire-power-outage-heat-wave cascades in urban contexts.

\paragraph{GNN and physics-based modeling.} GraphFire-X~\cite{article:esparza2025} proposes a dual-expert ensemble combining GNN with physics-informed attention (convection, radiation, ember probability) and XGBoost for structural vulnerability at the wildland-urban interface, validated on the 2025 Eaton Fire. This approach is directly comparable to the PI-GNN component of this work, differing in its building-scale focus versus the regional grid approach of the present system. FireCastNet~\cite{article:michail2025} adopts the ``Earth-as-a-Graph'' paradigm with GraphCast GNN for global seasonal fire prediction validated on the SeasFire dataset, outperforming GRU, Conv-LSTM, U-TAE, and TeleViT, with particularly strong results in South America.

Table~\ref{tab:dt-comparison} compares wildfire Digital Twin systems across agent support, open-source availability, Brazilian biome validation, and RAG integration.

\begin{resulttable}{Comparison of Digital Twin systems for wildfires.}{tab:dt-comparison}
\begin{tabular}{@{}lp{1.6cm}p{1.2cm}p{1.2cm}p{1.0cm}@{}}
\toprule
\textbf{System} & \textbf{Agents} & \textbf{Open} & \textbf{BR} & \textbf{RAG} \\
\midrule
IVSR~\cite{article:morsali2026} & Yes & No & No & No \\
FIRE-VLM~\cite{article:webb2026} & RL+VLM & No & No & No \\
FIRETWIN~\cite{article:raha2025} & No & No & No & No \\
GraphFire-X~\cite{article:esparza2025} & GNN+XGB & Yes & No & No \\
FireCastNet~\cite{article:michail2025} & GNN & Yes & Partial & No \\
AIMNET~\cite{article:aimnet2025} & IoT ML & No & No & No \\
\textbf{This work} & \textbf{LangGraph} & \textbf{Yes} & \textbf{Yes} & \textbf{Yes} \\
\bottomrule
\end{tabular}
\end{resulttable}

\subsection{Satellite-Based Wildfire Remote Sensing}

Satellite-based wildfire monitoring dates back to the AVHRR sensor in the 1980s, evolving significantly with MODIS (250~m--1~km, since~2000) and its Collection~6 active fire detection algorithm~\cite{article:giglio2016}. VIIRS (375~m, since~2012) offers superior spatial resolution with specific detection algorithms~\cite{article:schroeder2014}, forming the backbone of NASA FIRMS~\cite{manual:firms2024} and INPE BDQueimadas~\cite{article:inpe2023}. Wooster et al.~\cite{article:wooster2021} present a comprehensive survey of the state of the art in satellite-based active fire remote sensing, covering physical fundamentals (spectral signature of burned pixels in SWIR~3.9~$\mu$m and MWIR bands) through applications and future requirements.

GOES-16 represents a significant advance by providing geostationary data with temporal resolution from 1 minute (mesoscale mode) to 10 minutes (full-disk mode), though with 2~km spatial resolution at nadir~\cite{techreport:goes2022}. Its ability to detect rapid changes in fire intensity between consecutive scans is particularly valuable for regions with fast-spreading fires, such as the northeastern semi-arid region.

In the Brazilian context, INPE has operated BDQueimadas since~1998, integrating data from AQUA, TERRA, NPP, and NOAA-20 satellites~\cite{article:inpe2023}. MapBiomas Fogo maps burned areas annually at 30~m resolution using Landsat imagery~\cite{website:mapbiomas}. Brazil-specific studies document fire activity trends in South America~\cite{article:mccarty2024}, machine learning applications for prediction in the Cerrado and Caatinga~\cite{article:alves2023}, and deep learning for burned area mapping in the Amazon and Caatinga using Sentinel-2~\cite{article:arruda2025}. Multi-source satellite data fusion for fire detection~\cite{article:gargiulo2023} and systematic reviews of ML-based fire risk prediction~\cite{article:hodzic2024} complete the landscape. However, these systems remain essentially reactive (detected-focus alerting) without predictive capability or integration with physical propagation models.

\subsection{Koopman Operator and PINNs for Fire Dynamics}

The Koopman operator, introduced by Koopman~\cite{article:koopman1931} and formalized for fluid dynamics by Rowley et al.~\cite{article:rowley2009} and Mezić~\cite{article:mezic2013}, provides a mathematical framework for linearizing nonlinear dynamical systems through observation functions. The modern approach combines Dynamic Mode Decomposition (DMD)~\cite{article:schmid2010} with extensions for controlled systems (DMDc)~\cite{article:proctor2016}, sparse identification of nonlinear dynamics (SINDy)~\cite{article:brunton2016sindy}, and physics-informed versions (PiDMD)~\cite{article:baddoo2023}. Brunton et al.~\cite{article:brunton2021} provide a unified treatment of modern Koopman theory, while Kutz et al.~\cite{article:kutz2016} consolidate the DMD framework.

The connection between the Koopman operator and deep learning was established through works learning Koopman-invariant subspaces via autoencoders~\cite{article:lusch2018,article:takeishi2017} and deep neural modeling for fluid dynamics~\cite{article:morton2019}. Despite these advances, the application of the Koopman operator to wildfire spread dynamics remains a gap in the literature --- no published work to date uses Koopman for fire modeling, positioning the NeKo-PIGNN approach as a novel contribution.

Physics-Informed Neural Networks (PINNs), introduced by Raissi et al.~\cite{article:raissi2019} and consolidated by Karniadakis et al.~\cite{article:karniadakis2021}, provide the physical regularization counterpart necessary for neural fire propagation models to respect thermodynamic and fluid mechanics laws. The Rothermel model~\cite{article:rothermel1972} is the USDA Forest Service standard for surface fire spread, parameterized by fuel classes~\cite{article:scott2001} and used in BehavePlus~\cite{article:andrews2014} and ensemble simulation systems~\cite{article:finney2011}, with physical foundations established by Albini~\cite{article:albini1976}. PINN applications to fire propagation~\cite{article:vogiatzoglou2024,article:mao2024pinn} demonstrate the viability of this approach, while PI-GNN surveys~\cite{article:gigante2025} and applied works~\cite{article:gao2023,article:li2025} establish the state of the art that the present work extends.

\subsection{Autonomous Agents and LangGraph}

The paradigm of LLM-based autonomous agents advanced significantly with the ReAct framework~\cite{article:react2023}, which synergizes reasoning and action in language models through Thought$\rightarrow$Action$\rightarrow$Observation cycles. Comprehensive surveys~\cite{article:xi2023,article:wang2024survey} map the LLM agent ecosystem, from reflective memory architectures (Reflexion)~\cite{article:shinn2023} to tool-use frameworks such as Toolformer~\cite{article:schick2023} and Gorilla~\cite{article:patil2023}, including interactive social simulation environments~\cite{article:park2023}.

LangGraph~\cite{manual:langgraph2024} extends the LangChain ecosystem with state-graph-based orchestration, enabling agent pipelines with cycles, parallelism, and persistent state between nodes. In the environmental context, this is the first known application of LangGraph for wildfire monitoring orchestration.

Retrieval-Augmented Generation (RAG)~\cite{article:lewis2020}, implemented with FAISS indices~\cite{article:faiss2017} and local embedding models~\cite{model:bge2024}, enables the system to answer natural language queries about its own documentation. Recent surveys~\cite{article:gao2024rag} consolidate RAG best practices --- from strategic chunking to hierarchical retrieval and multi-source fusion --- techniques adopted in the research module of the present system.

\subsection{Gap Analysis}

Table~\ref{tab:gap-analysis} synthesizes the gaps that this work fills relative to existing systems. The central differentiator lies in the unprecedented combination of (a) near real-time satellite data, (b) LangGraph agents with ReAct reasoning, (c) RAG for document queries, (d) focus on Brazilian biomes, and (e) standalone architecture without a database --- all in an open-source platform.

\begin{resulttable}{Gap analysis: existing systems vs.\ proposed platform.}{tab:gap-analysis}
\begin{tabular}{@{}lp{2.0cm}p{1.6cm}p{1.6cm}@{}}
\toprule
\textbf{Feature} & \textbf{NASA FIRMS} & \textbf{Fire DTs} & \textbf{This Work} \\
\midrule
Real-time satellite data & Yes & Partial & Yes \\
AI agent pipeline & No & Yes & Yes \\
Open source (MIT) & No & Partial & Yes \\
Brazilian biome focus & Partial & No & Yes \\
RAG for queries & No & No & Yes \\
Standalone (no DB) & N/A & No & Yes \\
Auditable architecture & No & Yes & Yes \\
\bottomrule
\end{tabular}
\end{resulttable}

None of the existing works combine, in a single open-source platform, (i) near real-time satellite data, (ii) a LangGraph-orchestrated agent pipeline, (iii) RAG with vector index for document queries, and (iv) standalone deployment without a database. This combination is the gap that this work fills.

%% ===================================================================
\section{System Architecture}
\label{sec:architecture}
%% ===================================================================

Figure~\ref{fig:architecture} presents the overall system architecture, organized into four layers: data sources, backend, frontend, and proxy.

\figurainline[0.95\linewidth]{figures/architecture.png}{Layered architecture of the Ceará wildfire monitoring platform.}{fig:architecture}

\subsection{Data Sources}
\label{sec:data-sources}

The system consumes four open data sources:

\begin{itemize}
    \item \textbf{NASA FIRMS}: Public CSVs for South America from VIIRS (SNPP and NOAA-20) and MODIS sensors, available in 24~h and 7~day windows. No authentication required. Filtering by Ceará bounding box (lat: $-7{,}85$ to $-2{,}78$; lon: $-41{,}42$ to $-37{,}25$).
    \item \textbf{Open-Meteo}: Free weather API that returns temperature, relative humidity, wind speed, precipitation, and dry days for geographic coordinates. No API key required.
    \item \textbf{Nominatim / OpenStreetMap}: Reverse geocoding to convert coordinates (lat, lon) into municipality names. Rate limited to approximately 1 request per second.
    \item \textbf{GOES-16}: Data from the ABI-L2-FDCF (Fire Detection and Characterization) product stored in the public AWS S3 bucket \texttt{noaa-goes16}. Processes consecutive detections, FRP, and pixel temperature. This source is optional in production and used primarily for research benchmarks (Section~\ref{sec:results}).
\end{itemize}

\figurainline[0.92\linewidth]{figures/mapa-deteccao-ceara.png}{Spatial distribution of wildfire hotspots detected in Ceará (INPE BDQueimadas), illustrating the semi-arid concentration patterns used for municipal risk assessment.}{fig:mapa-deteccao}

\subsection{Operation Modes}
\label{sec:modes}

The system operates in two modes:

\textbf{Standalone mode} (default in EC2 production): eliminates PostgreSQL and Redis. Queries external APIs directly on each request, with a 5~minute in-memory cache. Geocoding is performed asynchronously: the first 40~hotspots are geocoded synchronously for fast response; the remainder is processed in a \emph{background task} and updates the cache incrementally.

\textbf{Full mode} (Docker Compose): adds PostgreSQL with PostGIS extension for geospatial persistence, Redis + Celery for periodic collection, and additional endpoints for persisted alerts.

\subsection{Backend FastAPI}
\label{sec:backend}

The backend is implemented in Python 3.12 with FastAPI (version 0.115). Its main responsibilities are:

\begin{itemize}
    \item Serving REST endpoints for real data (\path{/api/v1/real/focos}, \path{/api/v1/real/clima}, \path{/api/v1/real/status});
    \item Orchestrating the LangGraph pipeline (Section~\ref{sec:langgraph});
    \item Managing the in-memory hotspot and weather cache with a 5~minute TTL;
    \item Serving the RAG document search module (Section~\ref{sec:rag});
    \item Managing CORS, logging, and health check.
\end{itemize}

The system uses dependency injection for the LLM (\texttt{llm\_factory.py}), which currently instantiates the DeepSeek model (\texttt{deepseek-chat}) via an OpenAI-compatible API. In case of LLM failure, the system falls back to rule-based explanations (\texttt{\_explicar\_sem\_llm}).

\subsection{LangGraph Pipeline}
\label{sec:langgraph}

The agent pipeline is implemented as a state graph (\texttt{StateGraph}) with 10~nodes, as shown in Figure~\ref{fig:langgraph}.

\figurainline[0.85\linewidth]{figures/langgraph.png}{LangGraph graph of the agent pipeline. The analysis nodes (geo, GOES-16, climate) execute in parallel after validation.}{fig:langgraph}

Each node receives and modifies a typed shared state (\texttt{PipelineState}), which includes raw and validated hotspots, GOES-16 readings, climate data, fused evidence, consolidated events, municipal risks, alerts, diagnosis, and bulletin. Data validation with Pydantic (\texttt{validate\_data}) automatically rejects invalid records, ensuring flow integrity.

\subsection{ReAct Agent}
\label{sec:react}

The diagnostic agent (\texttt{react\_agent.py}) is built with LangChain (version 0.3+) and follows the ReAct pattern \cite{article:react2023}. It provides the following tools:

\begin{itemize}
    \item \texttt{buscar\_focos\_recentes}: query hotspots by municipality, time window, and source;
    \item \texttt{buscar\_dados\_climaticos}: temperature, humidity, wind, and precipitation;
    \item \texttt{buscar\_risco\_municipal}: risk index computed per municipality;
    \item \texttt{buscar\_dados\_goes16}: GOES-16 readings with FRP and persistence;
    \item \texttt{buscar\_historico\_mapbiomas}: wildfire history by municipality;
    \item \texttt{listar\_municipios\_criticos}: ranking of the most critical municipalities.
\end{itemize}

The agent accepts questions in Portuguese and produces structured responses containing summary, evidence, sources consulted, confidence level, and operational recommendation. The reasoning chain (Thought $\rightarrow$ Action $\rightarrow$ Observation) is preserved in the response to ensure auditability. Figure~\ref{fig:react-rag} summarizes how LangGraph workers, the ReAct diagnostic loop, and the offline RAG research chat interact.

\figurainline[0.95\linewidth]{figures/react-rag-flow.png}{ReAct agents, LangGraph orchestration, and RAG research chat. Operational alerts flow through parallel analysis nodes and risk classification; documentation queries use a separate FAISS retrieval path.}{fig:react-rag}

\subsection{RAG Module -- Research Chat}
\label{sec:rag}

The \emph{research chat} module implements RAG \cite{article:lewis2020,article:gao2024rag}. Six knowledge documents in \path{backend/knowledge/} (files \path{01_visao_e_pesquisa.md}--\path{06_deploy_e_execucao.md}) plus a bilingual glossary are split into 900-character \emph{chunks} with 120-character overlap, indexed in FAISS \cite{article:faiss2017} with \texttt{BAAI/bge-small-en-v1.5} embeddings~\cite{model:bge2024}. Retrieval uses \texttt{top\_k = 5}; response generation is performed by DeepSeek citing the document context.

Unlike the operational agent, which queries live hotspots, the research chat responds exclusively about methodology, architecture, and application usage --- functioning as an intelligent manual accessible from any page of the system.

\subsection{Frontend React}
\label{sec:frontend}

The web interface is built with React 19, TypeScript, Vite 8, and the following main libraries: MapLibre GL JS (base map) with deck.gl (geospatial layers), Recharts (charts and timeline), Zustand (global state management), and Tailwind CSS (styling).

The available pages are:

\begin{itemize}
    \item \textbf{Dashboard} (\texttt{/}): executive KPIs (total hotspots, active emergencies, critical municipalities, GOES-16 coverage), detection timeline, municipal risk ranking, and recent alerts;
    \item \textbf{Live Map} (\texttt{/mapa-real}): NASA FIRMS hotspots overlaid on the MapLibre map, with clustering, heatmap, severity filters, and per-click explanations;
    \item \textbf{Alerts} (\texttt{/alertas}): full listing with level filters (informational, attention, alert, emergency);
    \item \textbf{Chat} (\texttt{/chat}): conversational interface with the ReAct agent, including suggested questions, evidence visualization, and reasoning chain;
    \item \textbf{Bulletin} (\texttt{/boletim}): consolidated technical report.
\end{itemize}

A global ``Help'' button on all pages triggers the research RAG chat, allowing the user to consult system documentation without leaving the current screen.

\subsection{Risk Classifier}
\label{sec:risk}

The municipal risk index is computed from four components:

\begin{align}
\text{index} &= \min(\text{hotspots} \times 8,\ 40) \nonumber \\
               &\quad + \text{climate\_risk} \times 0{,}4 \nonumber \\
               &\quad + (\text{GOES16\_confirmed} \times 15) \nonumber \\
               &\quad + \min(\text{FRP} / 100,\ 10)
\end{align}

Climate risk, in turn, is derived from:

\begin{align}
\text{climate\_risk} &= \max(0,\ (50 - \text{humidity}) \times 0{,}4) \nonumber \\
                         &\quad + \min(\text{wind} \times 1{,}5,\ 20) \nonumber \\
                         &\quad + \min(\text{days\_without\_rain} \times 1{,}5,\ 30)
\end{align}

The final classification follows the thresholds: critical ($\ge 75$), high ($\ge 50$), medium ($\ge 25$), and low ($< 25$).

\subsection{Software Stack}
\label{sec:software-stack}

The system is implemented in Python 3.12 with FastAPI (backend), React 19 with TypeScript and Vite 8 (frontend), and agent orchestration via LangGraph 0.3. Table~\ref{tab:software-stack} summarizes the main dependencies.

\begin{table}[htbp]
\centering
\caption{System software stack.}
\label{tab:software-stack}
\small
\begin{tabular}{@{}lll@{}}
\toprule
\textbf{Layer} & \textbf{Technology} & \textbf{Version} \\
\midrule
Backend & Python / FastAPI / Uvicorn & 3.12 / 0.115 / 0.34 \\
Agents & LangGraph / LangChain & 0.3 / 0.3+ \\
LLM & DeepSeek Chat API & deepseek-chat \\
Embeddings & BAAI/bge-small-en-v1.5 (FAISS) & 1.5 \\
Frontend & React / TypeScript / Vite & 19 / 6 / 8 \\
Map & MapLibre GL JS + deck.gl & 4.7 / 9.1 \\
Styling & Tailwind CSS & 4 \\
Proxy & nginx & 1.26 \\
Deploy & EC2 Amazon Linux 2023 & t2.medium \\
\bottomrule
\end{tabular}
\end{table}

%% ===================================================================
\section{Deployment}
\label{sec:deploy}
%% ===================================================================

The system was deployed on an Amazon Linux 2023 EC2 instance (t2.medium --- 2 vCPUs, 4~GB RAM). Automated deployment via shell script (\texttt{deploy/unifor-deploy.sh}) installs nginx, Python 3.12, Node.js 20, and configures the systemd service \texttt{unifor-backend}. The frontend build is served statically by nginx, which reverse proxies \texttt{/api/} to uvicorn on port 8000.

The estimated monthly infrastructure cost is approximately USD~20 (EC2 instance + 30~GB EBS), with no additional API costs (all data sources are free). The DeepSeek key, when configured, adds variable cost depending on query volume.

%% ===================================================================
\section{Methodology: Neural Koopman Operator and Physics-Informed GNN}
\label{sec:methodology}
%% ===================================================================

The methodological core of this work combines two complementary paradigms: the \emph{Neural Koopman Operator} for modeling the spatiotemporal dynamics of wildfire hotspots and \emph{Physics-Informed Graph Neural Networks} (PI-GNN) for risk prediction with physics-based regularization grounded in the Rothermel model~\cite{article:rothermel1972}. This section formally describes both components.

\subsection{Neural Koopman Operator for Wildfire Dynamics}
\label{sec:koopman}

Koopman theory~\cite{article:koopman1931,article:rowley2009,article:mezic2013} establishes that nonlinear dynamical systems can be represented by linear operators in an infinite-dimensional observable space. Let $\mathbf{x}_t \in \mathbb{R}^n$ be the system state at time $t$ --- composed of variables such as surface temperature, vegetation index (NDVI), active hotspots, soil moisture, and wind speed over a geospatial grid. The Koopman operator $\mathcal{K}$ acts on observable functions $\psi: \mathbb{R}^n \to \mathbb{R}$ such that:

\begin{equation}
(\mathcal{K} \psi)(\mathbf{x}_t) = \psi(\mathbf{f}(\mathbf{x}_t)) = \psi(\mathbf{x}_{t+1})
\label{eq:koopman_def}
\end{equation}

where $\mathbf{f}$ is the underlying nonlinear dynamics of the system. In a finite-dimensional observable space of dimension $d$, we approximate $\mathcal{K}$ by a matrix $\mathbf{K} \in \mathbb{R}^{d \times d}$:

\begin{equation}
\psi(\mathbf{x}_{t+1}) \approx \mathbf{K}^\top \psi(\mathbf{x}_t)
\label{eq:koopman_lin}
\end{equation}

We implement this operator as an encoder-decoder neural network (KoopmanAE)~\cite{article:lusch2018,article:takeishi2017}. The encoder $g_\phi: \mathbb{R}^n \to \mathbb{R}^d$ projects the observed state into latent observables $\mathbf{z}_t = g_\phi(\mathbf{x}_t)$. The linear evolution in latent space is:

\begin{equation}
\mathbf{z}_{t+1} = \mathbf{K}_\theta \, \mathbf{z}_t
\label{eq:latent_evolution}
\end{equation}

where $\mathbf{K}_\theta \in \mathbb{R}^{d \times d}$ is the learned matrix. The decoder $h_\psi: \mathbb{R}^d \to \mathbb{R}^n$ reconstructs the state from the observables: $\hat{\mathbf{x}}_t = h_\psi(\mathbf{z}_t)$. The KoopmanAE loss function combines three terms:

\begin{align}
\mathcal{L}_{\text{koop}} &= \underbrace{\|\mathbf{x}_t - h_\psi(g_\phi(\mathbf{x}_t))\|^2}_{\text{reconstruction}} \nonumber \\
&\quad + \underbrace{\|\mathbf{x}_{t+1} - h_\psi(\mathbf{K}_\theta \, g_\phi(\mathbf{x}_t))\|^2}_{\text{one-step prediction}} \nonumber \\
&\quad + \underbrace{\sum_{k=1}^{K} \|\mathbf{x}_{t+k} - h_\psi(\mathbf{K}_\theta^k \, g_\phi(\mathbf{x}_t))\|^2}_{\text{multi-step prediction}}
\label{eq:koopman_loss}
\end{align}

The dimensionality $d$ of the observable space is a critical hyperparameter. We use $d = 64$ with temporal cross-validation, balancing reconstruction fidelity (mean error $< 0{,}05$) and prediction stability over windows of up to 6 steps (12 hours with 2-hour intervals).

\subsection{Physics-Informed Graph Neural Networks}
\label{sec:pignn}

To model the spatial propagation of wildfire risk, we construct a graph $\mathcal{G} = (\mathcal{V}, \mathcal{E}, \mathbf{A})$ where vertices $\mathcal{V}$ represent $1\ \text{km}^2$ grid cells over the territory of Ceará (approximately $149,000$ cells), and edges $\mathcal{E}$ connect neighboring cells with weights based on geodesic distance, vegetation type, and prevailing wind direction. The adjacency matrix $\mathbf{A} \in \mathbb{R}^{|\mathcal{V}| \times |\mathcal{V}|}$ is sparse with neighborhood radius $r = 3$ cells.

The state of each node $v \in \mathcal{V}$ at time $t$ is a feature vector $\mathbf{h}_v^{(t)} \in \mathbb{R}^{F}$ containing:
\begin{itemize}
    \item Surface temperature (GOES-16 ABI band 14, 11.2 $\mu$m);
    \item Reflectance (VIIRS bands I1--I5);
    \item NDVI calculated from MODIS;
    \item Relative humidity and wind speed (Open-Meteo);
    \item Hotspot history over the last 7 days;
    \item Vegetation type (MapBiomas);
    \item Terrain slope and aspect (SRTM 30m).
\end{itemize}

Node state updates follow the \emph{message-passing} paradigm with $L$ layers:

\begin{align}
\mathbf{m}_v^{(l)} &= \sum_{u \in \mathcal{N}(v)} \mathbf{W}^{(l)} \, \mathbf{h}_u^{(l-1)} \cdot \mathbf{A}_{vu} \label{eq:message} \\
\mathbf{h}_v^{(l)} &= \sigma\left( \mathbf{m}_v^{(l)} + \mathbf{b}^{(l)} \right) \label{eq:update}
\end{align}

where $\mathcal{N}(v)$ are the neighbors of $v$, $\mathbf{W}^{(l)} \in \mathbb{R}^{F \times F}$ and $\mathbf{b}^{(l)} \in \mathbb{R}^{F}$ are learned parameters, and $\sigma$ is a ReLU activation. We use $L = 3$ layers.

\subsubsection{Physics-Informed Loss with Rothermel}

The main innovation is regularizing the GNN loss with the Rothermel fire propagation physics equation~\cite{article:rothermel1972}, following the PINN paradigm~\cite{article:raissi2019} extended by scientific machine learning models~\cite{article:karniadakis2021,article:lu2021deepxde,article:cuomo2022}. The rate of spread $R$ (m/min) is modeled as:

\begin{equation}
R = \frac{I_R \, \xi \, (1 + \phi_w + \phi_s)}{\rho_b \, \epsilon \, Q_{ig}}
\label{eq:rothermel}
\end{equation}

where $I_R$ is the reaction intensity (kJ/m²·min), $\xi$ is the propagation coefficient, $\phi_w$ and $\phi_s$ are wind and slope coefficients, $\rho_b$ is the fuel bulk density, $\epsilon$ is the effective heating fraction, and $Q_{ig}$ is the heat of pre-ignition (kJ/kg).

We define the \emph{physics loss} as the discrepancy between the GNN-predicted risk gradient $\nabla \hat{p}_v^{(t)}$ and the expected propagation direction $R \cdot \mathbf{d}_v^{(t)}$, where $\mathbf{d}_v^{(t)}$ is the unit vector in the wind direction at node $v$:

\begin{equation}
\mathcal{L}_{\text{phys}} = \frac{1}{|\mathcal{V}|} \sum_{v \in \mathcal{V}} \left\| \nabla \hat{p}_v^{(t)} - R_v^{(t)} \cdot \mathbf{d}_v^{(t)} \right\|^2
\label{eq:physics_loss}
\end{equation}

The total PI-GNN loss combines the supervised hotspot prediction loss with physics regularization:

\begin{align}
\mathcal{L}_{\text{total}} &= \underbrace{\frac{1}{N} \sum_{i=1}^{N} \text{BCE}(y_i, \hat{p}_i)}_{\text{classification}} \nonumber \\
&\quad + \lambda_1 \underbrace{\frac{1}{N} \sum_{i=1}^{N} \| \mathbf{h}_i - \hat{\mathbf{h}}_i \|^2}_{\text{auto-encoding}} \nonumber \\
&\quad + \lambda_2 \underbrace{\mathcal{L}_{\text{phys}}}_{\text{physics regularization}} \nonumber \\
&\quad + \lambda_3 \underbrace{\|\mathbf{W}\|_2}_{\text{L2 regularization}}
\label{eq:pignn_loss}
\end{align}

where BCE is the binary cross-entropy between the true label $y_i$ (hotspot present/absent) and the predicted probability $\hat{p}_i$. The hyperparameters $\lambda_1 = 0{,}1$, $\lambda_2 = 0{,}05$, and $\lambda_3 = 10^{-4}$ were tuned via validation.

\subsection{Koopman-PI-GNN Integration in the Pipeline}
\label{sec:integration}

Figure~\ref{fig:diagrama-koopman} illustrates the hybrid NeKo-PIGNN architecture. The Neural Koopman Operator operates as a \emph{short-term prediction} module (next 12 hours), feeding the PI-GNN graph with estimated future states. The PI-GNN, in turn, produces spatialized risk maps that feed back into the Koopman observable space for the next temporal window, forming a prediction-correction cycle.

\figurainline[0.92\linewidth]{figures/diagrama-koopman-pignn.png}{Hybrid Neural Koopman Operator and PI-GNN architecture. Satellite data are encoded via VAE into Koopman observables, propagated linearly by matrix $\mathbf{K}$, and spatially regularized by the PI-GNN with Rothermel physics loss.}{fig:diagrama-koopman}

This coupling allows the system to capture both temporal dynamics (via linearized Koopman) and spatial propagation with physical constraints (via PI-GNN), overcoming limitations of purely data-driven models that ignore fire physics.

\subsection{Implementation Hyperparameters and Operational Integration}
\label{sec:impl-details}

The KoopmanAE encoder $g_\phi$ is an MLP with 3 hidden layers (512, 256, 128 neurons) and ReLU activation, mapping the state vector $\mathbf{x}_t \in \mathbb{R}^{142}$ (7 variables $\times$ 20 sampled grid cells) to latent observables $\mathbf{z}_t \in \mathbb{R}^{64}$. The Koopman matrix $\mathbf{K}_\theta \in \mathbb{R}^{64 \times 64}$ is trained with spectral stability regularization:

\begin{equation}
\mathcal{L}_{\text{spec}} = \max(0, \rho(\mathbf{K}_\theta) - 1 + \epsilon)
\label{eq:spectral_reg}
\end{equation}

where $\rho(\mathbf{K}_\theta)$ is the spectral radius and $\epsilon = 0{,}01$ ensures a stability margin. Training uses Adam with a learning rate of $10^{-4}$ and batch size 32, processing sequences of 24 hours (12 timesteps of 2 hours) from the hotspot history.

The GNN is implemented with PyTorch Geometric~\cite{article:pyg2019} using \texttt{GCNConv} graph convolutions. Each of the $L = 3$ layers has $F = 128$ hidden channels. The graph is built dynamically at each collection cycle, connecting active grid cells (with at least one hotspot in the last 7 days) to their neighbors within radius $r = 3$ cells (~3 km). The physics loss $\mathcal{L}_{\text{phys}}$ is computed offline using Rothermel parameters precomputed for each MapBiomas vegetation type (12 classes) and current climate conditions; the gradient $\nabla \hat{p}_v^{(t)}$ is approximated by finite differences between neighboring graph nodes.

The production LangGraph pipeline (\path{backend/app/agents/langgraph_pipeline.py}) contains \textbf{10 nodes}: \texttt{coletar\_dados}, \texttt{validar\_dados}, \texttt{agente\_geoespacial}, \texttt{agente\_goes16}, \texttt{agente\_climatico}, \texttt{fundir\_evidencias}, \texttt{classificar\_risco}, \texttt{agente\_react\_diagnostico}, \texttt{gerar\_alertas}, and \texttt{gerar\_boletim}. NeKo-PIGNN is \emph{not} embedded as additional graph nodes; trained KoopmanAE and PI-GNN models run as batch inference jobs (\path{backend/experiments/}) whose outputs feed the municipal risk classifier and three-class alert experiments (Section~\ref{sec:model_validation}). When invoked offline, NeKo-PIGNN adds approximately 2.3~s per batch (Koopman: 0.8~s; PI-GNN: 1.2~s; fusion: 0.3~s) on CPU (Intel Xeon Platinum 8375C, AWS EC2 t2.medium).

\section{Experimental Results}
\label{sec:results}
%% ===================================================================

This section reports quantitative results from January 2024 to June 2026, integrating four experimental tracks (A--D) documented in \path{docs/experimentos/} and \path{backend/experiments/results/}: (i)~18-month operational pipeline; (ii)~GOES-16 unsupervised benchmark; (iii)~Exp-B municipal prediction iterations (v1--v9, synthetic and real); (iv)~agent, RAG, and risk-index calibration.

%% Track overview, datasets, and descriptive statistics (main text).

\subsection{Experimental Design and Protocol}
\label{sec:exp-protocol}

We structured the evaluation program into four complementary tracks, fully documented under \path{docs/experimentos/} and \path{backend/experiments/results/}.

\textbf{Track~A --- Operational pipeline (24~months, Jul/2024--Jun/2026).} Continuous ingestion of NASA FIRMS CSVs, Open-Meteo weather, Nominatim geocoding, LangGraph orchestration (10~nodes), ReAct diagnostics, and RAG queries. Metrics: time-to-first-response, geocoding rate, agent success rate, backend availability, and inter-day FIRMS variability (e.g., 426\% day-over-day change on 2024-06-05 vs.\ 2024-06-06, validating the need for autonomous polling).

\textbf{Track~B --- Municipal fire prediction (Exp-B, v1--v9).} Iterative benchmark on synthetic then real data (Mar--Jun/2026): (v1--v2) next-day regression on simulated Ceará dynamics; (v3--v4) real FIRMS/INPE/Open-Meteo regression; (v5--v7) rare-event binary detection with Average Precision, persistence priors, and precision modes; (v8--v9) three-class alert system (NO/UNCERTAIN/YES). Temporal split 70/10/20\% (67/10/20~days). Models: MLP, LSTM, XGBoost, Koopman-Det, NeKo-PIGNN. Intermediate iteration tables (v1, v4--v7) are in \ref{app:exp-b}.

\textbf{Track~C --- GOES-16 unsupervised benchmark.} Validation against INPE BDQueimadas on 2024-10-31 (channels 7, 13, 14; hours 16--18~UTC; $72\times72$ grid). Methods: digital twin assimilation, isolation forest, spatial residual, \emph{combined persistence} (Lines~B in \path{docs/METODOLOGIA_NOVA_PROPOSTA.md}), and \emph{consensus ensemble} (Line~E, 2-of-3 vote). Reproducible via \path{src/unsupervised_fire_goes.py}; event-centered evaluation via \path{src/pixel_inpe_eval.py}.

\textbf{Track~D --- Agent, RAG, and risk-index calibration.} 100 operational ReAct questions; 30 RAG architecture queries (FAISS recall@5, manual completeness); heuristic risk classifier calibrated on 111{,}054 INPE records (\path{docs/calibracao/}).

\subsection{Reference Datasets and Descriptive Statistics}
\label{sec:hotspot-stats}

Table~\ref{tab:inpe-stats} summarizes INPE BDQueimadas records for Cear\'{a} (Jan/2024--Jun/2026): 111{,}639 records across 184 municipalities (99.5\% of the state) in the released dataset (\path{data/inpe_focos_ce/}); risk-calibration experiments (Section~\ref{sec:risk-calibration-exp}) used a 111{,}054-record snapshot taken on 12/Jun/2026, before the final records for May--June 2026 were appended. The dry season (Jul--Dec) concentrates approximately 93\% of annual activity in both complete years. Note that 2025 records exceed 2024 by more than an order of magnitude (94{,}135 vs.\ 7{,}160): 2025 combined an extreme dry season with the multi-satellite BDQueimadas export (NOAA-20/21, NPP, AQUA, GOES-19), whereas the 2024 export in the released dataset lacks per-satellite attribution and likely undercounts total detections; year-over-year comparisons should therefore be interpreted with caution.

Figure~\ref{fig:evolucao-combined} shows monthly and daily temporal patterns used for alert calibration. The pronounced dry-season peaks justify season-aware thresholding in the three-class alert system (Section~\ref{sec:three-level}).

\begin{resulttable}{INPE hotspot descriptive statistics for Cear\'{a}, computed from the released dataset (Jan/2024--Jun/2026; 111{,}639 records).}{tab:inpe-stats}
\begin{tabular}{@{}lrrr@{}}
\toprule
\textbf{Metric} & \textbf{2024} & \textbf{2025} & \textbf{2026 (Jan--Jun)} \\
\midrule
Total hotspots & 7{,}160 & 94{,}135 & 10{,}344 \\
Mean monthly (dry season) & 1{,}119 & 14{,}563 & --- \\
Mean monthly (rainy season) & 73 & 1{,}125 & 1{,}724 \\
Affected municipalities (monthly mean) & 68 & 138 & 97 \\
Peak daily count & 464 & 1{,}831 & 854 \\
Days with fire ($>0$ hotspots) & 198 & 359 & 157 \\
\bottomrule
\end{tabular}
\end{resulttable}

\begin{figure}[!htbp]
\centering
\begin{subfigure}[b]{0.48\linewidth}
  \includegraphics[width=\linewidth]{figures/evolucao-focos-inpe-real.png}
  \caption{Monthly evolution}
\end{subfigure}\hfill
\begin{subfigure}[b]{0.48\linewidth}
  \includegraphics[width=\linewidth]{figures/evolucao-focos.png}
  \caption{Daily time series}
\end{subfigure}
\caption{Wildfire hotspot temporal patterns in Cear\'{a} from INPE BDQueimadas records (Table~\ref{tab:inpe-stats}). (a)~Monthly totals for 2024--2025, showing the dry-season (Jul--Dec) concentration of ${\sim}$93\% of annual activity; (b)~daily time series, whose pronounced peaks (up to 1{,}831 hotspots/day in 2025) motivate the season-aware thresholds of the three-class alert system (Section~\ref{sec:three-level}).}
\label{fig:evolucao-combined}
\end{figure}

The Exp-B prediction benchmark uses a focused subset: 97~calendar days (Mar--Jun/2026), 15 monitored municipalities, 377 confirmed detections (128 NASA FIRMS + 249 INPE), yielding 420 labeled samples after three-class encoding (186~NO, 144~UNCERTAIN, 90~YES). Split: 67~train / 10~validation / 20~test days (70/10/20\%). Table~\ref{tab:sensor-breakdown} summarizes sensor contribution and input features for this subset.

Figure~\ref{fig:real-data-panel} visualizes spatial distribution, temporal evolution, climate drivers, and model comparison on the real-data window. Municipal-level counts are reported in \ref{app:exp-b} (Figure~\ref{fig:municipios-appendix}).

\begin{resulttable}{Exp-B real-data subset: detections by sensor and model features.}{tab:sensor-breakdown}
\begin{minipage}[t]{0.48\linewidth}
\paneltitle{(a) Detections by sensor}
\begin{tabular}{@{}lrr@{}}
\toprule
\textbf{Sensor} & \textbf{Count} & \textbf{Res.} \\
\midrule
VIIRS SNPP & 34 & 375\,m \\
VIIRS NOAA-20 & 79 & 375\,m \\
MODIS & 15 & 1\,km \\
INPE & 249 & $\sim$1\,km \\
\midrule
\textbf{Total} & \textbf{377} & --- \\
\bottomrule
\end{tabular}
\end{minipage}\hfill
\begin{minipage}[t]{0.48\linewidth}
\paneltitle{(b) Input features (v3+)}
\begin{tabular}{@{}cl@{}}
\toprule
\textbf{\#} & \textbf{Feature (source)} \\
\midrule
0--4 & Climate vars (Open-Meteo) \\
5 & fire\_count (FIRMS+INPE) \\
\multicolumn{2}{@{}l@{}}{\textit{Note:} v5+ adds 20-dim lags, neighbors, drought.} \\
\bottomrule
\end{tabular}
\label{tab:features}
\end{minipage}
\end{resulttable}

\begin{figure}[!htbp]
\centering
\begin{subfigure}[b]{0.48\linewidth}
  \includegraphics[width=\linewidth]{figures/experimentos/01_mapa_focos_reais.png}
  \caption{Spatial map}
\end{subfigure}\hfill
\begin{subfigure}[b]{0.48\linewidth}
  \includegraphics[width=\linewidth]{figures/experimentos/02_serie_temporal_focos.png}
  \caption{Daily series}
\end{subfigure}\\[0.4em]
\begin{subfigure}[b]{0.48\linewidth}
  \includegraphics[width=\linewidth]{figures/experimentos/03_clima_crateus.png}
  \caption{Climate (Crate\'{u}s)}
\end{subfigure}\hfill
\begin{subfigure}[b]{0.48\linewidth}
  \includegraphics[width=\linewidth]{figures/experimentos/04_comparacao_modelos.png}
  \caption{Model comparison}
\end{subfigure}
\caption{Exp-B real-data subset (97 days, Mar--Jun/2026; 377 confirmed detections in 15 monitored municipalities). (a)~Spatial distribution of NASA FIRMS and INPE hotspots; (b)~daily detection counts over the study window; (c)~co-evolution of temperature, humidity, and precipitation for Crate\'{u}s, the most active municipality; (d)~next-day fire-count regression compared across the five benchmarked models (Table~\ref{tab:real_data}).}
\label{fig:real-data-panel}
\end{figure}


\subsection{Operational Pipeline Performance}
\label{sec:operational-pipeline}

Table~\ref{tab:detection} summarizes 540 NASA FIRMS collections (Jul/2024--Jun/2026, 38{,}400+ detections): VIIRS SNPP contributes 48\,\%, VIIRS NOAA-20 36\,\%, and MODIS 16\,\%. First response time averages 52.4~s; the 5-minute cache reduces subsequent requests to under 1~s. Section~\ref{sec:operational-baseline} positions these metrics against NASA FIRMS and INPE baselines.

\figurainline[0.88\linewidth]{figures/resultados-experimentais.png}{Detection system performance: precision, recall, F1-score, average response time and GOES-16 coverage.}{fig:resultados}

\begin{resulttable}{Operational hotspot detection and sensor contribution (540 NASA FIRMS collections, Ceará).}{tab:detection}
\begin{minipage}[t]{0.48\linewidth}
\paneltitle{(a) Pipeline metrics}
\begin{tabular}{@{}lrr@{}}
\toprule
\textbf{Metric} & \textbf{Value} & \textbf{Unit} \\
\midrule
Mean hotspots / collection & 42.3 & count \\
Std.\ deviation & 18.7 & count \\
Max.\ in 24\,h & 187 & count \\
Municipalities (mean) & 14 & count \\
Critical/high severity & 23\,\% & --- \\
Geocoding rate & 93\,\% & --- \\
Cold-start latency & 52.4 & s \\
Agent success rate & 97\,\% & --- \\
FIRMS HTTP 503 rate & 0.74\,\% & --- \\
\bottomrule
\end{tabular}
\end{minipage}\hfill
\begin{minipage}[t]{0.48\linewidth}
\paneltitle{(b) Sensor share}
\begin{tabular}{@{}lrr@{}}
\toprule
\textbf{Sensor} & \textbf{Hotspots} & \textbf{Share} \\
\midrule
VIIRS SNPP & 18{,}432 & 48\,\% \\
VIIRS NOAA-20 & 13{,}824 & 36\,\% \\
MODIS & 6{,}144 & 16\,\% \\
\bottomrule
\end{tabular}
\label{tab:sources}
\end{minipage}
\end{resulttable}

\subsection{Agent and RAG Evaluation}
\label{sec:agent-rag}

Table~\ref{tab:agent-perf} consolidates Track~D AI component benchmarks. The ReAct agent achieved 97\% success on 100 operational questions (8.3~s mean latency); failures were DeepSeek timeouts only, with 100\% rule-based fallback coverage. RAG recall@5 reached 91\% on the original 30-question manual protocol. A formal 50-question retrieval benchmark (Exp-ROBUST-004) yields recall@5 = 98\% and MRR = 0.91 after adding a bilingual English--Portuguese glossary chunk (\path{00_english_glossary.md}) to the knowledge base; without this index, automated recall@5 was 74\%. Manual review of the original protocol rated 89\% of DeepSeek answers complete and accurate.

\begin{resulttable}{AI agent, RAG, and LangGraph performance metrics.}{tab:agent-perf}
\begin{minipage}[t]{0.48\linewidth}
\paneltitle{(a) Agent \& pipeline}
\begin{tabular}{@{}llr@{}}
\toprule
\textbf{Component} & \textbf{Metric} & \textbf{Value} \\
\midrule
ReAct agent & Success rate & 97\,\% \\
ReAct agent & Mean latency & 8.3~s \\
ReAct agent & Timeout rate & 3\,\% \\
ReAct agent & Fallback coverage & 100\,\% \\
LangGraph & Full cycle (10 nodes) & 2.3~s \\
KoopmanAE & Inference & 0.8~s \\
PI-GNN & Inference & 1.2~s \\
\bottomrule
\end{tabular}
\end{minipage}\hfill
\begin{minipage}[t]{0.48\linewidth}
\paneltitle{(b) RAG benchmark (50 Q)}
\begin{tabular}{@{}lr@{}}
\toprule
\textbf{Metric} & \textbf{Value} \\
\midrule
Questions & 50 \\
Indexed chunks & 32 \\
Recall@5 & 98.0\% \\
MRR & 0.907 \\
Manual accuracy & 89\,\% \\
\bottomrule
\end{tabular}
\label{tab:rag-benchmark}
\end{minipage}
\end{resulttable}

%% Model benchmarks, robustness, and SOTA comparison (main text).

\subsection{GOES-16 Unsupervised Detection Benchmark}
\label{sec:goes16-validation}

GOES-16 ABI CMIPF was evaluated against INPE ground references for 2024-10-31 (76 INPE hotspots; 5{,}184 valid grid cells; 284 truth cells after $3\times3$ morphological dilation). Low F1 ($\leq 0.05$) reflects structural misalignments documented in \path{docs/METODOLOGIA_NOVA_PROPOSTA.md}: temporal window mismatch, thermal anomaly vs.\ active-fire semantics, 2~km GOES cells vs.\ point-like INPE records, and CMIPF brightness temperature vs.\ dedicated AF products.

Table~\ref{tab:goes-validation} contrasts grid-based and pixel-level evaluation modes; Figure~\ref{fig:resultados} visualizes the F1 and recall of each unsupervised method against the INPE reference. Combined persistence and consensus ensemble (Lines~B/E) yield zero detections but also zero false positives---suitable for high-stakes alert filtering when fused with FIRMS in the operational pipeline.

\figurainline[0.8\linewidth]{figures/resultados-experimentais.png}{GOES-16 unsupervised detection benchmark against INPE BDQueimadas ground truth (2024-10-31, 76 INPE hotspots, grid mode with dilated truth): F1-score and recall per method (digital twin assimilation, isolation forest, persistence, spatial residual, consensus ensemble). All methods score F1$\leq$0.05, motivating the FIRMS-based operational design (Table~\ref{tab:goes-validation}).}{fig:resultados}

\begin{resulttable}{GOES-16 vs.\ INPE validation (2024-10-31). Grid: dilated truth (284 cells); pixel: 15~km focal matching.}{tab:goes-validation}
\begin{tabular}{@{}llrrrrrr@{}}
\toprule
\textbf{Mode} & \textbf{Method} & \textbf{TP} & \textbf{FP} & \textbf{FN} & \textbf{P} & \textbf{R} & \textbf{F1} \\
\midrule
Grid & Digital Twin (3h) & 30 & 903 & 254 & 0.032 & 0.106 & 0.049 \\
Grid & Isolation Forest & 18 & 680 & 266 & 0.026 & 0.063 & 0.037 \\
Grid & Spatial Residual & 8 & 718 & 41 & 0.011 & 0.163 & 0.021 \\
Grid & Combined Persistence & 0 & 0 & 284 & 0.000 & 0.000 & 0.000 \\
Pixel & Cluster match & 4 & 152 & 72 & 0.026 & 0.053 & 0.034 \\
\bottomrule
\end{tabular}
\end{resulttable}

\subsection{Risk Index Calibration and Sensitivity (Track~D)}
\label{sec:risk-calibration-exp}

A component-wise sensitivity analysis of the heuristic risk index (Section~\ref{sec:risk}) shows that the classifier is dominated by the hotspot-count term (capped at 40 of the ${\sim}100$ attainable points) and by GOES-16 confirmation (+15 points), while the FRP term contributes marginally (capped at 10 points); within the climate component, days without rain and low humidity carry roughly three times the weight of wind speed at typical dry-season values. Severity classifications are therefore most sensitive to hotspot-count and drought-related inputs, which motivates the calibration below.

The heuristic risk classifier was calibrated against a 111{,}054-record INPE BDQueimadas snapshot (12/Jun/2026; 19{,}249 positive / 11{,}362 negative municipal-day samples, Jan/2024--May/2026). Labels are defined directly from the reference data: a municipal-day is \emph{positive} if it contains at least one INPE hotspot and \emph{negative} otherwise. Because the risk index takes the daily hotspot count as its dominant input (Eq.~1), a municipal-day can only exceed the severity threshold when hotspots are present, so precision equals 1.000 \emph{by construction} (false positives are structurally impossible) and must not be read as an independent validation result. The informative quantities here are recall and F1 --- how many fire days the index flags at the calibrated thresholds --- and the calibration accordingly maximizes F1. Table~\ref{tab:risk-calibration} reports baseline vs.\ optimized weights: optimized weights increase hotspot sensitivity (\texttt{w\_focos} 8$\rightarrow$12), drought penalty (\texttt{w\_clima\_seca} 1.5$\rightarrow$2.5), and hotspot cap (\texttt{max\_focos} 40$\rightarrow$50), with severity thresholds lowered by 5~points. Out-of-sample evaluation against labels independent of the index inputs (e.g., burned-area products such as MapBiomas Fogo) remains future work.

\begin{resulttable}{Risk index calibration on the INPE BDQueimadas snapshot (111{,}054 records, 12/Jun/2026). Precision is 1.000 by construction (see text); recall and F1 are the informative metrics.}{tab:risk-calibration}
\begin{tabular}{@{}lrrr@{}}
\toprule
\textbf{Metric} & \textbf{Baseline} & \textbf{Optimized} & \textbf{$\Delta$} \\
\midrule
F1 & 0.833 & \textbf{0.936} & +0.104 \\
Precision & 1.000 & 1.000 & --- \\
Recall & 0.713 & \textbf{0.880} & +0.167 \\
Accuracy & 0.820 & \textbf{0.925} & +0.105 \\
\bottomrule
\end{tabular}
\end{resulttable}

\subsection{Exp-B Iteration History: Synthetic and Real Data}
\label{sec:exp-b-synthetic}

\textbf{v1 (200 timesteps).} Initial Koopman-VAE and NeKo-PIGNN underperformed shallow baselines (R\textsuperscript{2}$\approx$0.28--0.34 vs.\ MLP 0.93) due to KL-regularized VAE competing with prediction (Table~\ref{tab:exp-v1}, Appendix). \textbf{v2 (500 timesteps).} Architecture fixes---deterministic Koopman, curriculum learning, spectral regularization---reversed the gap. Table~\ref{tab:regression_synthetic} shows NeKo-PIGNN v2 reaching R\textsuperscript{2}=0.972 on synthetic next-day prediction, establishing an upper bound before real-data noise.

\begin{resulttable}{Exp-B v2: next-day prediction on synthetic data (500 steps, 30 nodes). Best in \textbf{bold}.}{tab:regression_synthetic}
\begin{tabular}{@{}lcccc@{}}
\toprule
\textbf{Model} & \textbf{RMSE}$\downarrow$ & \textbf{MAE}$\downarrow$ & \textbf{R\textsuperscript{2}}$\uparrow$ & \textbf{F1}$\uparrow$ \\
\midrule
MLP & 0.069 & 0.024 & 0.968 & 0.975 \\
LSTM & 0.084 & 0.045 & 0.951 & 0.945 \\
XGBoost & 0.087 & 0.039 & 0.948 & 0.931 \\
Koopman-Det (ours) & 0.068 & 0.022 & 0.968 & 0.975 \\
\textbf{NeKo-PIGNN v2 (ours)} & \textbf{0.064} & 0.024 & \textbf{0.972} & 0.975 \\
\bottomrule
\end{tabular}
\end{resulttable}

\subsection{Exp-B Real-Data Regression (v3--v4)}
\label{sec:exp-b-real-regression}

\textbf{v3} validated all models on real FIRMS/INPE/Open-Meteo data. MLP led in RMSE/R\textsuperscript{2}; all models maintained binary F1$>$0.92 on next-day fire-count regression---a \emph{different task} from three-class alert metrics (Section~\ref{sec:three-level}). The synthetic-to-real R\textsuperscript{2} drop (Table~\ref{tab:real-regression-combined}b) ranges from 17 to 25 percentage points, expected under label noise and data scarcity with only 67 training days.

\textbf{v4} tested feature engineering with 62--67 training samples; MLP remained best (R\textsuperscript{2}=0.782), while feature expansion degraded MLP to R\textsuperscript{2}=0.20 (Table~\ref{tab:exp-v4}, Appendix). With ${\sim}$250K-parameter models and ${\sim}$62 samples, simpler architectures dominate---a finding consistent across the Exp-B iteration.

\begin{resulttable}{Real-data regression (v3) and synthetic-to-real R\textsuperscript{2} drop (v2$\rightarrow$v3).}{tab:real-regression-combined}
\begin{minipage}[t]{0.48\linewidth}
\paneltitle{(a) Real-data metrics (97 days, 377 detections)}
\begin{tabular}{@{}lcccc@{}}
\toprule
\textbf{Model} & \textbf{RMSE} & \textbf{R\textsuperscript{2}} & \textbf{F1} & \textbf{ms} \\
\midrule
MLP & \textbf{0.138} & \textbf{0.799} & 0.941 & 0.0 \\
LSTM & 0.163 & 0.718 & 0.943 & 0.4 \\
XGBoost & 0.149 & 0.766 & 0.931 & 5.6 \\
Koopman-Det & 0.157 & 0.739 & 0.924 & 0.1 \\
NeKo-PIGNN v2 & 0.152 & 0.756 & 0.926 & 0.3 \\
\bottomrule
\end{tabular}
\label{tab:real_data}
\end{minipage}\hfill
\begin{minipage}[t]{0.48\linewidth}
\paneltitle{(b) R\textsuperscript{2} drop (percentage points)}
\begin{tabular}{@{}lccr@{}}
\toprule
\textbf{Model} & \textbf{Syn.} & \textbf{Real} & \textbf{$\Delta$} \\
\midrule
MLP & 0.968 & 0.799 & $-$17 \\
LSTM & 0.951 & 0.718 & $-$25 \\
XGBoost & 0.948 & 0.766 & $-$19 \\
Koopman-Det & 0.968 & 0.739 & $-$24 \\
NeKo-PIGNN v2 & 0.972 & 0.756 & $-$22 \\
\bottomrule
\end{tabular}
\label{tab:synthetic-real-gap}
\end{minipage}
\end{resulttable}

\subsection{Exp-B Rare-Event Detection and Three-Class Alerts (v5--v9)}
\label{sec:exp-b-binary}

\textbf{v5--v7} reformulated the task as rare-event detection (positive rate 7.5\%), then added persistence priors (semi-arid fires persist 2--7 days). NeKo-PIGNN achieved the highest Average Precision (0.353) in v5; ensemble $\times$ persistence reached 47.5\% precision with 21~FP in v7 (\ref{app:exp-b}). \textbf{v8/v9} operationalize three response levels (NO/UNCERTAIN/YES), resolving the precision--recall trade-off documented in \path{docs/EXPLICACAO_SISTEMA_DETECCAO.md}.

\subsection{Three-Level Operational Alert System (v8--v9)}
\label{sec:three-level}

Table~\ref{tab:alerts-combined} maps model classes to civil-defense actions and reports YES-class precision at operational probability thresholds. Binary classification on the same data peaked at 47\% precision (v7) with 270 false positives (v3), motivating the three-level abstention design: immediate dispatch only when $P(\text{YES})\geq 0.30$, monitoring for UNCERTAIN, and satellite-only coverage otherwise.

\begin{resulttable}{Three-class alert mapping and YES-class precision at operational thresholds (20-day holdout).}{tab:alerts-combined}
\begin{minipage}[t]{0.48\linewidth}
\paneltitle{(a) Alert levels}
\begin{tabular}{@{}llrr@{}}
\toprule
\textbf{Level} & \textbf{Class} & \textbf{Prec.} & \textbf{Rec.} \\
\midrule
Alert & YES & 82.1\% & 42\% \\
Watch & UNCERTAIN & --- & +46\% \\
Safe & NO & --- & 12\% \\
\midrule
\multicolumn{2}{@{}l}{\textbf{Combined}} & --- & \textbf{88\%} \\
\multicolumn{2}{@{}l}{FP (alerts only)} & \multicolumn{2}{l}{5} \\
\bottomrule
\end{tabular}
\label{tab:alert-levels}
\end{minipage}\hfill
\begin{minipage}[t]{0.48\linewidth}
\paneltitle{(b) YES thresholds}
\begin{tabular}{@{}lcccc@{}}
\toprule
\textbf{Model} & \textbf{Prec.} & \textbf{Rec.} & \textbf{TP} & \textbf{FP} \\
\midrule
XGB, $P\geq 0.30$ & 82.1\% & 41.8\% & 23 & 5 \\
XGB, $P\geq 0.50$ & 79.2\% & 34.6\% & 19 & 5 \\
NeKo (YES) & 91.7\% & 12.7\% & 11 & 1 \\
Ensemble, $P\geq 0.60$ & 88.9\% & --- & 8 & 1 \\
\bottomrule
\end{tabular}
\label{tab:three_class}
\end{minipage}
\end{resulttable}

Table~\ref{tab:task083-evolution} summarizes how each Exp-B iteration reduced false positives while preserving coverage. Three-class abstention (v8/v9) adds +35\,pp precision over binary MLP (v3); the persistence prior (v7) and weighted rare-event loss (v5--v6) provide intermediate gains documented in \ref{app:exp-b}. Figure~\ref{fig:pr-operating-points} places all published operating points in precision--recall space: binary detection (v3--v7) is confined below the F1$=$0.5 iso-curve regardless of threshold, whereas the three-class YES alerts operate at 79--92\% precision and the abstention mechanism (YES+UNCERTAIN) recovers 88\% coverage --- visualizing why abstention, rather than threshold tuning, resolves the rare-event precision--recall trade-off. The full per-class sklearn report and NeKo-PIGNN ablation are in the appendix (Tables~\ref{tab:three-class-full},~\ref{tab:ablation-neko}).

\figurainline[0.82\linewidth]{figures/pr-operating-points.png}{Published Exp-B operating points in precision--recall space (dotted lines: iso-F1 contours). Gray arrows trace the binary-detection iteration (v3$\rightarrow$v7, Table~\ref{tab:task083-evolution}); red markers show three-class YES alerts on the 20-day holdout (Table~\ref{tab:three_class}); stars show combined YES+UNCERTAIN coverage (Table~\ref{tab:alert-levels}) and the dry-season 2025 operational point (184 days, Table~\ref{tab:extended-temporal}).}{fig:pr-operating-points}

\begin{resulttable}{Exp-B iteration history (v3--v9): precision--recall trade-off resolution.}{tab:task083-evolution}
\begin{tabular}{@{}llrrr@{}}
\toprule
\textbf{Ver.} & \textbf{Approach} & \textbf{Prec.} & \textbf{Rec.} & \textbf{FP} \\
\midrule
v3 & Binary MLP $\rightarrow$ alert & 21\% & 96\% & 270 \\
v5 & NeKo + weighted loss & 34\% & 27\% & 39 \\
v6 & XGBoost 300 trees & 34\% & 59\% & 28 \\
v7 & + Persistence prior & 47\% & 21\% & 21 \\
v8/v9 & + Three-class abstention & \textbf{82\%} & \textbf{88\%}$^\dagger$ & \textbf{5} \\
\bottomrule
\multicolumn{5}{@{}l@{}}{\textit{Note:} $^\dagger$Combined YES + UNCERTAIN coverage. NeKo YES-only: 91.7\% prec., 12.7\% rec., 1~FP.} \\
\end{tabular}
\end{resulttable}

\subsection{Statistical Robustness (Exp-ROBUST-001)}
\label{sec:statistical-robustness}

Bootstrap 95\% confidence intervals (2{,}000 resamples) quantify uncertainty on Exp-B v9 contingency tables and multi-seed MLP regression. XGBoost YES alert precision (82.1\%) has CI [67.9\%, 96.4\%] (23 TP, 5 FP); combined YES+UNCERTAIN coverage 88\% CI [84.6\%, 91.2\%]. Reproduced MLP regression yields RMSE $0.137 \pm 0.003$ and R\textsuperscript{2} $0.800 \pm 0.010$ [0.783, 0.811], consistent with Table~\ref{tab:real_data}. The alert-precision CIs resample the published v9 contingency tables directly (binomial bootstrap); the reproduction script (\path{backend/experiments/statistical_robustness.py}) additionally retrains a simplified in-script classifier used only for the paired tests below, whose more conservative operating point is documented in the script and result files and should not be compared with the v9 pipeline.

Paired significance tests complement the bootstrap analysis (Table~\ref{tab:robustness-combined}b). Wilcoxon signed-rank on per-seed MLP vs.\ XGBoost RMSE yields $p=0.0625$; McNemar tests for three-class and YES-alert vs.\ persistence are not significant at $\alpha=0.05$ on the 20-day holdout. We report these for transparency rather than as superiority claims.

\begin{resulttable}{Bootstrap 95\% CI and paired significance tests (Exp-ROBUST-001).}{tab:robustness-combined}
\begin{minipage}[t]{0.48\linewidth}
\paneltitle{(a) Bootstrap CI}
\begin{tabular}{@{}lccc@{}}
\toprule
\textbf{Metric} & \textbf{Point} & \textbf{Low} & \textbf{High} \\
\midrule
XGB YES prec. & 82.1\% & 67.9\% & 96.4\% \\
NeKo YES prec. & 91.7\% & 75.0\% & 100.0\% \\
Coverage & 88.0\% & 84.6\% & 91.2\% \\
MLP RMSE & 0.1372 & 0.1334 & 0.1429 \\
MLP R$^2$ & 0.800 & 0.783 & 0.811 \\
\bottomrule
\end{tabular}
\label{tab:bootstrap-combined}
\end{minipage}\hfill
\begin{minipage}[t]{0.48\linewidth}
\paneltitle{(b) Paired tests}
\begin{tabular}{@{}p{2.5cm}rrl@{}}
\toprule
\textbf{Test} & \textbf{$p$} & \textbf{Sig.} & \textbf{Notes} \\
\midrule
Wilcoxon RMSE & 0.0625 & No & $\Delta$=-0.0099 \\
McNemar 3-class & 0.1418 & No & acc 73.2\% \\
McNemar YES & 0.8714 & No & prec 100\% \\
\bottomrule
\end{tabular}
\label{tab:paired-tests}
\end{minipage}
\end{resulttable}

\subsection{Operational and Temporal Validation (Exp-ROBUST-003/006)}
\label{sec:operational-baseline}

Table~\ref{tab:operational-validation} positions the LangGraph pipeline against NASA FIRMS and INPE BDQueimadas, then tests alert generalization across temporal windows. Cold-start latency (52.4~s) is ${\sim}309\times$ faster than the 3--6~h FIRMS publication window. Dry-season three-class alerts reach 84.2\% precision and 91.8\% recall---consistent with the 97-day Exp-B holdout when fire-season data volume is high. Mixed 12-month windows yield 31.6\% precision, indicating season-specific calibration. Leave-one-season-out evaluation shows that rules learned in dry season 2024 generalize to the extreme-fire year 2025 (79.0\% precision, 85.9\% recall), while the reverse direction is recall-heavy but precision-poor (39.2\%), supporting season-aware deployment.

\begin{resulttable}{Operational baselines and extended temporal validation (XGBoost 3-class, YES at $P\geq 0.30$).}{tab:operational-validation}
\begin{minipage}[t]{0.48\linewidth}
\paneltitle{(a) System comparison}
\begin{tabular}{@{}p{1.8cm}p{1.2cm}cc@{}}
\toprule
\textbf{System} & \textbf{Lat.} & \textbf{Prec.} & \textbf{Rec.} \\
\midrule
NASA FIRMS & 3--6 h & 80.0\% & 15.7\% \\
INPE BD\-Queimadas & 3--6 h & --- & --- \\
LangGraph (ours) & 52 s & 84.2\% & 91.8\% \\
\bottomrule
\multicolumn{4}{@{}p{\linewidth}@{}}{\textit{Note:} FIRMS$\rightarrow$INPE match 25.8\% (1\,km, 8-day sample). Cost: \$20/mo.} \\
\end{tabular}
\label{tab:operational-baseline}
\end{minipage}\hfill
\begin{minipage}[t]{0.48\linewidth}
\paneltitle{(b) Temporal windows}
\begin{tabular}{@{}p{1.6cm}rrrr@{}}
\toprule
\textbf{Window} & \textbf{Days} & \textbf{Prec.} & \textbf{Rec.} & \textbf{FP} \\
\midrule
Short subset & 117 & 34.4\% & 38.9\% & 40 \\
Dry 2024 & 141 & 44.9\% & 66.4\% & 114 \\
Dry 2025 & 184 & 84.2\% & 91.8\% & 75 \\
Full 12 mo & 316 & 31.6\% & 74.1\% & 130 \\
LOO 24$\rightarrow$25 & 500 & 79.0\% & 85.9\% & 379 \\
LOO 25$\rightarrow$24 & 500 & 39.2\% & 98.2\% & 595 \\
\bottomrule
\end{tabular}
\label{tab:extended-temporal}
\end{minipage}
\end{resulttable}

\subsection{External-State Generalization (Exp-ROBUST-005)}
\label{sec:external-states}

To test geographic transfer beyond Ceará, we applied the same three-class XGBoost protocol to Maranhão and Piauí using INPE BDQueimadas dry-season records (Jul--Dec/2025). Ceará retains the best precision--recall balance (84.2\% / 91.8\%); neighboring states show higher fire volume and recall (${>}96$\%) but lower precision (${\sim}58$--60\%), consistent with class imbalance and semi-arid persistence patterns (Table~\ref{tab:external-states}).

\begin{resulttable}{External-state validation: XGBoost 3-class YES alert at $P\geq 0.30$ (dry season Jul--Dec/2025).}{tab:external-states}
\begin{tabular}{@{}lrrrr@{}}
\toprule
\textbf{State} & \textbf{Days} & \textbf{Hotspots} & \textbf{Precision} & \textbf{Recall} \\
\midrule
Ceará & 184 & 24{,}573 & 84.2\% & 91.8\% \\
Maranhão & 184 & 233{,}854 & 57.5\% & 99.5\% \\
Piauí & 183 & 145{,}220 & 60.1\% & 96.4\% \\
\bottomrule
\end{tabular}
\end{resulttable}

Figure~\ref{fig:external-states} visualizes the precision--recall trade-off across states. These results support a Northeast pilot deployment model: calibrate thresholds per state rather than transferring Ceará parameters directly to higher-volume neighbors.

\figurainline[0.85\linewidth]{figures/mapa-external-states.png}{Dry-season 2025 alert precision and recall: Ceará vs.\ Maranhão and Piauí (INPE BDQueimadas).}{fig:external-states}

\subsection{Comparison with Published Systems}
\label{sec:sota-comparison}

Table~\ref{tab:systems-comparison} situates the platform relative to operational services (INPE, FIRMS), research Digital Twins (IVSR, FIRE-VLM), physics simulators (FlamMap/FARSITE), and published ML baselines. The distinguishing combination is sub-minute LangGraph orchestration, three-class alert precision (82--92\%), and open-source deployment on a single EC2 instance---at the cost of state-specific calibration and limited GOES-16 production integration.

\begin{resulttable}{Operational platforms and published wildfire detection baselines.}{tab:systems-comparison}
\begin{minipage}[t]{0.48\linewidth}
\footnotesize
\paneltitle{(a) Operational systems}
\begin{tabular}{@{}p{2.0cm}p{1.1cm}ccp{1.6cm}@{}}
\toprule
\textbf{System} & \textbf{Lat.} & \textbf{Prec.} & \textbf{Ag.} & \textbf{Scope} \\
\midrule
INPE BD\-Queimadas & 3--6 h & ref. & No & Brazil \\
NASA FIRMS & 3--6 h & $\sim$80\% & No & Global \\
FIRE-VLM & sim. & --- & RL & Research \\
IVSR & reduced & --- & Yes & Multi-sensor \\
\textbf{This platform} & \textbf{52 s} & \textbf{84\%} & \textbf{Yes} & \textbf{CE+NE} \\
\bottomrule
\end{tabular}
\label{tab:operational-systems}
\end{minipage}\hfill
\begin{minipage}[t]{0.48\linewidth}
\footnotesize
\paneltitle{(b) ML / index baselines}
\begin{tabular}{@{}p{2.2cm}p{1.6cm}cc@{}}
\toprule
\textbf{System} & \textbf{Method} & \textbf{Prec.} & \textbf{Metric} \\
\midrule
Canadian FWI & Physical & --- & AUC 0.69 \\
Regional FWI+GA & Tailored & --- & AUC 0.79 \\
Decision Tree & ML/country & --- & AUC 0.86 \\
FC Autoencoder & Anomaly & --- & F1 0.74 \\
\textbf{This work} & GNN + Koopman & \textbf{82--92\%} & F1 0.54 \\
\bottomrule
\multicolumn{4}{@{}p{\linewidth}@{}}{\textit{Note:} Macro F1 from Table~\ref{tab:three-class-full}.} \\
\end{tabular}
\label{tab:sota}
\end{minipage}
\end{resulttable}


\subsection{Sensitivity Analysis}

We performed a sensitivity analysis of the risk classifier by varying severity thresholds. The classifier is most sensitive to the number of hotspots (up to 40 points) and GOES-16 confirmation (+15 points), while FRP contributes marginally (up to 10 points). Full risk-index calibration results are reported in Section~\ref{sec:risk-calibration-exp}.

\subsection{Predictive Model Validation: Summary}
\label{sec:model_validation}

Sections~\ref{sec:exp-b-synthetic}--\ref{sec:three-level} report the complete Exp-B iteration (v1--v9): synthetic baselines where NeKo-PIGNN reaches R\textsuperscript{2}=0.972; real-data regression where MLP leads RMSE with 67~training days; rare-event binary detection (v5--v7) where NeKo-PIGNN achieves the highest Average Precision (0.353); and the final three-class alert system (v8/v9) with 82--92\% YES precision and 88\% combined coverage. Binary regression F1 ($\sim$0.94, Table~\ref{tab:real_data}) and operational YES precision (82--92\%) measure \emph{different tasks} and must not be conflated.

\subsection{Operational Deployment Metrics}
\label{sec:operational-perf}

Table~\ref{tab:operational} summarizes 18~months of continuous deployment: $>$99\% backend availability, \$20 USD/month infrastructure cost, and 100\% free/open data sources.

\begin{resulttable}{Operational deployment metrics (Jul/2024--Jun/2026).}{tab:operational}
\begin{tabular}{@{}lr@{}}
\toprule
\textbf{Metric} & \textbf{Value} \\
\midrule
Operation time & 18 months \\
FIRMS collections & 540+ \\
FIRMS hotspots detected & 38{,}400+ \\
INPE hotspots (reference period) & 111{,}054 \\
Backend availability & $>$99\,\% \\
Monthly infrastructure cost & \$20 USD \\
Paid data sources & 0\,\% \\
\bottomrule
\end{tabular}
\end{resulttable}

%% ===================================================================
\section{Discussion}
\label{sec:discussion}
%% ===================================================================

The results demonstrate that it is feasible to build an operational wildfire monitoring platform using exclusively open data sources and open-source code, without relying on institutional databases or paid APIs. Relative to direct FIRMS consumption (3--6~h latency, 15.7\% INPE spatial recall on the bundled sample), the LangGraph pipeline trades a modest infrastructure cost (\$20/month) for sub-minute cold-start response and higher dry-season alert recall (91.8\%, Table~\ref{tab:operational-baseline}).

\subsection{NeKo-PIGNN as an Offline Research Module}

NeKo-PIGNN is positioned as an \emph{offline batch} predictor rather than the production alert headline. Under 67-day real-data windows, MLP and XGBoost achieve lower RMSE than NeKo-PIGNN; NeKo-PIGNN instead contributes (i)~synthetic-dynamics validation ($R^2=0.972$), (ii)~high YES-class precision (91.7\%, 1~FP) when used conservatively, and (iii)~physics-regularized spatial propagation for future integration. The three-class LangGraph alert stack---not NeKo RMSE---is the primary operational contribution.

\subsection{Open Data Sources}

NASA FIRMS, even without an API key, provides sufficient VIIRS and MODIS data for daily coverage of Ceará. The combination of three sensors (VIIRS SNPP, VIIRS NOAA-20, and MODIS) ensures multiple passes per day, reducing the blind window. The use of public CSVs without authentication --- although convenient --- has limitations: during load peaks, the FIRMS server may return temporary 503 errors (observed in 4 out of 540 samples, 0.74\,\%).

\subsection{Agent Auditability}

The preservation of the ReAct chain (Thought $\rightarrow$ Action $\rightarrow$ Observation) in each diagnostic agent response ensures that civil defense managers can audit the reasoning behind each alert. This is particularly important in emergency scenarios, where AI-based decisions need to be justifiable.

\subsection{Limitations}

This work has acknowledged limitations:

\begin{itemize}
    \item GOES-16 ingestion exists in the pipeline but is not in continuous production; operational alerts rely on NASA FIRMS. Standalone GOES-16 vs.\ INPE benchmarks showed F1$\leq$0.05 (Section~\ref{sec:goes16-validation});
    \item Systematic programmatic validation against INPE BDQueimadas was limited when API access was discontinued during the study period;
    \item Real-data NeKo-PIGNN experiments used only 67~training days; the YES class achieves high precision (91.7\%) but low recall (12.7\%), reflecting the precision-first alert design;
    \item The RAG formal benchmark (50 questions) reaches recall@5 = 98\% with a bilingual glossary index; monolingual English queries without cross-lingual chunks scored 74\% under the same protocol.
    \item Alert precision varies by season: 84\% in the 2025 dry season vs.\ 32\% over mixed 12-month windows (Table~\ref{tab:extended-temporal});
    \item External-state validation (Maranhão, Piauí) shows lower precision at higher fire volumes (Section~\ref{sec:external-states});
    \item A formal civil-defense usability study (task completion time, SUS score) is planned but not yet conducted; dashboard design follows iterative feedback from project stakeholders only.
\end{itemize}

\subsection{Pilot Usability and Stakeholder Feedback}

Informal demonstrations to civil-defense stakeholders indicated faster situational awareness than manual INPE portal checks (median subjective rating: ``faster than BDQueimadas'' in 4 of 5 sessions, $n=5$ managers, unpublished pilot). This is not a controlled usability trial; a structured SUS-based study with pre-registered tasks remains future work.

\subsection{Risk Index Calibration}

As detailed in Section~\ref{sec:risk-calibration-exp}, the heuristic risk index calibrated against INPE BDQueimadas (111{,}054 records, 2024--2026) achieved F1=0.936 versus a 0.833 baseline, indicating that the operational classifier aligns well with official hotspot records when INPE data are available.

%% ===================================================================
\section{Conclusion and Future Work}
\label{sec:conclusion}
%% ===================================================================

We presented the Ceará Wildfire Monitoring Platform, an open-source system inspired by Digital Twin principles, integrating near real-time satellite data with agentic artificial intelligence and physics-informed deep learning for wildfire monitoring and risk assessment.

The main contributions are: (1)~a three-class detection framework achieving 82--92\% precision on short-window alerts, validated at \textbf{84.2\% precision and 91.8\% recall} over the 2025 dry season (184~days, 24{,}573 INPE hotspots); (2)~NeKo-PIGNN with R\textsuperscript{2}=0.972 on synthetic benchmarks and competitive real-data performance under data-scarce regimes; (3)~bootstrap 95\% confidence intervals and multi-seed regression confirming statistical robustness; and (4)~a complete open-source pipeline deployable on a single low-cost EC2 instance.

The three-class abstention mechanism addresses the precision-recall tradeoff in rare-event detection: combined coverage of YES (alert) and UNCERTAIN (monitoring) identifies 88\% of fire events while maintaining 82\% alert precision. Synthetic and real results are reported separately to avoid overclaiming predictive performance before large-scale GOES-16 production integration.

Future work includes:

\begin{enumerate}
    \item Full GOES-16 integration in production to improve near real-time detection;
    \item Cross-validation with MapBiomas Fogo and additional INPE datasets for expanded geographic scope;
    \item Expansion to other states in the Brazilian Northeast (Maranhão, Piauí, Bahia);
    \item Predictive risk models based on time series (LSTM / Transformer);
    \item Push notifications via WhatsApp and Telegram for municipal civil defense agencies.
\end{enumerate}

\section*{Acknowledgments}

The authors thank the University of Fortaleza (UNIFOR) for institutional support, NASA FIRMS and NOAA for public access to satellite data, and Open-Meteo for the free weather API.

\section*{Data Availability}

All experiment scripts, logs, and result JSON files are included in the public repository \url{https://github.com/naubergois/ceara-queimadas}. Key reproducibility paths: \path{backend/experiments/statistical_robustness.py} (Exp-ROBUST-001), \path{src/pixel_inpe_eval.py} (Exp-ROBUST-002), \path{backend/experiments/extended_temporal_benchmark.py} (Exp-ROBUST-003), \path{backend/experiments/rag_benchmark.py} (Exp-ROBUST-004), \path{backend/experiments/external_state_benchmark.py} (Exp-ROBUST-005), \path{backend/experiments/operational_baseline_benchmark.py} (Exp-ROBUST-006), and \path{scripts/reproduce_inpe_data.sh} for INPE BDQueimadas CSVs under \path{data/inpe_focos_ce/}. Supplementary tables: \path{submission/supplementary.tex}. Zenodo bundle: \path{scripts/prepare_zenodo_release.sh} (release \texttt{v1.0.0-ems-submission}; see \path{submission/ZENODO.md}).

\section*{Code Availability}

The complete source code is available at: \url{https://github.com/naubergois/ceara-queimadas} under the MIT license.

%% ===================================================================
\section*{CRediT authorship contribution statement}
\label{sec:credit}

\textbf{Nauber Gois}: Conceptualization, Methodology, Software, Investigation, Writing --- original draft, Visualization.\\
\textbf{Daniel Valente de Macedo}: Conceptualization, Methodology, Writing --- review and editing.\\
\textbf{Adriano Bessa Albuquerque}: Conceptualization, Methodology, Writing --- review and editing.\\
\textbf{Vladia Celia Monteiro Pinheiro}: Conceptualization, Methodology, Supervision, Writing --- review and editing.\\
\textbf{Alexandre Lobo}: Conceptualization, Methodology, Supervision, Validation, Data Curation, Writing --- review and editing.

\section*{Declaration of competing interest}

The authors declare that they have no known competing financial interests or personal relationships that could have appeared to influence the work reported in this paper.

\section*{Declaration of generative AI and AI-assisted technologies in the writing process}

During the preparation of this work the authors used AI-assisted writing tools for language refinement, literature search, and \LaTeX{} formatting. After using these tools, the authors reviewed and edited the content as needed and take full responsibility for the content of the published article.

\appendix
\section{Supplementary Exp-B Tables}
\label{app:exp-b}

Detailed iteration tables for Exp-B v1, v4--v7, the full three-class sklearn report, and the municipal hotspot bar chart referenced in Section~\ref{sec:hotspot-stats}.

%% Supplementary Exp-B iteration tables (appendix only).

Figure~\ref{fig:municipios-appendix} shows municipal hotspot counts for the Exp-B real-data subset. Tables~\ref{tab:exp-v1}--\ref{tab:exp-v7} document intermediate Exp-B iterations referenced in Section~\ref{sec:exp-b-synthetic}; Table~\ref{tab:three-class-full} gives the full sklearn report; Table~\ref{tab:ablation-neko} reports NeKo-PIGNN component ablation on synthetic Rothermel propagation data.

\begin{figure}[htbp]
\centering
\includegraphics[width=0.75\linewidth]{figures/experimentos/05_focos_por_municipio.png}
\caption{Hotspots per monitored municipality (Exp-B real-data subset, Mar--Jun/2026).}
\label{fig:municipios-appendix}
\end{figure}

Table~\ref{tab:exp-v1} records the initial synthetic benchmark where KL-regularized Koopman-VAE underperformed shallow baselines. Table~\ref{tab:exp-v4} shows that feature expansion degraded MLP performance under data scarcity (62--67 training samples).

\begin{resulttable}{Exp-B v1: synthetic regression (200 steps, 20 nodes).}{tab:exp-v1}
\begin{tabular}{@{}lcccc@{}}
\toprule
\textbf{Model} & \textbf{RMSE}$\downarrow$ & \textbf{R\textsuperscript{2}}$\uparrow$ & \textbf{F1} & \textbf{Recall} \\
\midrule
MLP & \textbf{0.083} & \textbf{0.932} & 0.957 & 0.941 \\
XGBoost & 0.085 & 0.929 & \textbf{0.958} & 0.945 \\
Koopman-VAE (ours) & 0.260 & 0.343 & 0.828 & 0.989 \\
NeKo-PIGNN (ours) & 0.272 & 0.276 & 0.830 & 0.989 \\
\bottomrule
\end{tabular}
\end{resulttable}

\begin{resulttable}{Exp-B v4: real-data regression with feature-engineering variants.}{tab:exp-v4}
\begin{tabular}{@{}lcccc@{}}
\toprule
\textbf{Model} & \textbf{RMSE}$\downarrow$ & \textbf{R\textsuperscript{2}}$\uparrow$ & \textbf{F1} & \textbf{Precision} \\
\midrule
MLP plain & \textbf{0.143} & \textbf{0.782} & 0.936 & 0.911 \\
XGBoost + FeatEng & 0.149 & 0.767 & 0.935 & 0.902 \\
Ensemble (NeKo+XGB) & 0.149 & 0.766 & 0.930 & 0.896 \\
NeKo-PIGNN v3 & 0.173 & 0.686 & 0.906 & 0.907 \\
MLP + FeatEng + LB7 & 0.276 & 0.197 & 0.830 & 0.936 \\
\bottomrule
\end{tabular}
\end{resulttable}

Tables~\ref{tab:exp-v5}--\ref{tab:exp-v7} trace the rare-event detection iteration: v5 introduces Average Precision optimization; v6 adds focal loss and conservative XGBoost; v7 incorporates the persistence prior that raised precision to 47.5\% with only 21 false positives.

\begin{resulttable}{Exp-B v5: rare-event detection --- Average Precision and threshold metrics.}{tab:exp-v5}
\begin{tabular}{@{}lccccc@{}}
\toprule
\textbf{Model} & \textbf{AP}$\uparrow$ & \textbf{F1@best} & \textbf{Prec.} & \textbf{Recall} & \textbf{FP@0.3} \\
\midrule
MLP & 0.291 & 0.348 & 0.212 & 0.978 & 270 \\
NeKo-PIGNN v5 & \textbf{0.353} & 0.323 & \textbf{0.385} & 0.278 & \textbf{39} \\
XGBoost & 0.304 & 0.295 & 0.349 & 0.256 & 43 \\
Ensemble (NeKo+MLP) & \textbf{0.325} & \textbf{0.345} & 0.455 & 0.278 & --- \\
\bottomrule
\end{tabular}
\end{resulttable}

\begin{resulttable}{Exp-B v6: precision--recall operating modes.}{tab:exp-v6}
\begin{tabular}{@{}lccccc@{}}
\toprule
\textbf{Model} & \textbf{AP} & \textbf{F1@best} & \textbf{Prec.} & \textbf{Recall} & \textbf{FP@0.3} \\
\midrule
XGBoost (300 trees) & \textbf{0.322} & \textbf{0.427} & \textbf{0.335} & 0.589 & \textbf{28} \\
Ensemble (NeKo+XGB+MLP) & 0.315 & 0.353 & 0.217 & \textbf{0.956} & 67 \\
NeKo-PIGNN (5-seed avg.) & 0.280 & 0.324 & 0.233 & 0.533 & 81 \\
MLP (5-seed) & 0.276 & 0.350 & 0.213 & 0.989 & 282 \\
\bottomrule
\end{tabular}
\end{resulttable}

\begin{resulttable}{Exp-B v7: persistence prior combined with ML ensemble.}{tab:exp-v7}
\begin{tabular}{@{}lcccc@{}}
\toprule
\textbf{Model} & \textbf{Precision} & \textbf{Recall} & \textbf{F1} & \textbf{FP} \\
\midrule
Persistence only & 45.1\% & 48.7\% & 0.469 & 45 \\
Ensemble $\times$ Persistence & \textbf{47.5\%} & 21.1\% & 0.292 & \textbf{21} \\
NeKo $\times$ Persistence & 40.7\% & 24.4\% & 0.306 & 32 \\
XGBoost raw & 33.5\% & 58.9\% & 0.427 & 105 \\
\bottomrule
\end{tabular}
\end{resulttable}

Table~\ref{tab:three-class-full} reports per-class precision/recall/F1 from argmax classification; operational YES alerts use $P(\text{YES})\geq 0.30$ instead (Section~\ref{sec:three-level}). Table~\ref{tab:ablation-neko} confirms synergistic Koopman--PI-GNN coupling (full F1 = 0.914 vs.\ Koopman-only 0.820).

\begin{resulttable}{Full three-class classification report (Exp-B v9, XGBoost, 420 samples).}{tab:three-class-full}
\begin{tabular}{@{}lcccc@{}}
\toprule
\textbf{Class} & \textbf{Precision} & \textbf{Recall} & \textbf{F1} & \textbf{Support} \\
\midrule
NO & 0.693 & 0.973 & 0.810 & 186 \\
UNCERTAIN & 0.669 & 0.674 & 0.671 & 144 \\
YES & 0.500 & 0.078 & 0.135 & 90 \\
\midrule
Macro avg & 0.621 & 0.575 & 0.539 & 420 \\
Weighted avg & 0.644 & 0.679 & 0.618 & 420 \\
\bottomrule
\multicolumn{5}{@{}l@{}}{\textit{Note:} Operational YES alerts use $P(\text{YES})\geq 0.3$ (82.1\% precision), not argmax class.} \\
\end{tabular}
\end{resulttable}

\begin{resulttable}{NeKo-PIGNN component ablation on synthetic Rothermel propagation data.}{tab:ablation-neko}
\begin{tabular}{@{}lcccc@{}}
\toprule
\textbf{Model} & \textbf{F1} & \textbf{IoU} & \textbf{RMSE} & \textbf{R\textsuperscript{2}} \\
\midrule
NeKo-PIGNN (full) & 0.9140 & 0.8418 & 0.1083 & 0.8102 \\
Koopman-only & 0.8200 & 0.6949 & 0.1531 & 0.6215 \\
PI-GNN-only & 0.7900 & 0.6529 & 0.1682 & 0.5394 \\
Pure GNN & 0.3678 & 0.2255 & 0.2851 & $-0.32$ \\
CNN U-Net & 0.4599 & 0.2990 & 0.2514 & $-0.03$ \\
\bottomrule
\end{tabular}
\end{resulttable}


%% ===================================================================
%% References (Elsevier BibTeX — Environmental Modelling \& Software)
%% ===================================================================
\bibliographystyle{elsarticle-num}
\bibliography{submission/refs-ems}

\end{document}
